\chapter*{Resumen}
\addcontentsline{toc}{chapter}{Resumen}

\textbf{Resumen (ES):}

El presente informe documenta el diseño, implementación y validación del Sistema de Gestión de Recursos Operativos, Administrativos y de Seguridad para Cooperativas de Transporte (SGROAS), aplicación web de tres capas desarrollada como Proyecto Fin de Curso en la Universidad Técnica Estatal de Quevedo. SGROAS gestiona recursos operativos, administrativos y de seguridad mediante arquitectura basada en Angular 17+ (frontend), Spring Boot 3.2.x con Java 21 LTS (backend), PostgreSQL 16 y Redis 7. Ofrece autenticación stateless mediante JWT en cookie HttpOnly, CRUD completo para conductores, usuarios, rutas, vehículos, asignaciones e incidentes, seis procedimientos almacenados/funciones SQL y reportes. La verificación se realizó mediante 150 pruebas JUnit 5 con cobertura JaCoCo del 85,5\% líneas / 54,3\% ramas en el núcleo (67,5\% líneas full con \texttt{abd} excluido del \texttt{jacoco:check}), pruebas de carga k6 (3 corridas, 50 VUs, 30\,s) con media de medias 23,01\,ms IC95\% [-7,88; 53,91], p95 máximo 173,13\,ms (umbral 200\,ms, 0\% errores) y análisis de seguridad OWASP con ZAP baseline. Los requisitos se especificaron conforme a ISO/IEC/IEEE 29148:2018 y se verificaron las 15 características del INCOSE v4. El SUS obtuvo media 63,0 (DT 13,88, IC95\% [53,07; 72,93], n=10). El sistema se versiona como v1.0.0 con DOI Zenodo.

\textbf{Palabras clave:} ingeniería de requisitos, SRS, Spring Boot, Angular, JWT, PostgreSQL, OWASP, ISO 29148, INCOSE, proyecto fin de curso.

\medskip

\textbf{Abstract (EN):}

This report documents the design, implementation and validation of the Operations, Administrative and Security Resource Management System for Transport Cooperatives (SGROAS), a three-tier web application developed as Final Course Project at the State Technical University of Quevedo. SGROAS manages operational, administrative and security resources through an architecture based on Angular 17+ (frontend), Spring Boot 3.2.x with Java 21 LTS (backend), PostgreSQL 16 and Redis 7. It provides stateless JWT authentication in HttpOnly cookies, full CRUD for drivers, users, routes, vehicles, assignments and incidents, six stored procedures/SQL functions and reports. Verification included 150 JUnit 5 tests with JaCoCo 85.5\% lines / 54.3\% branches on the core (67.5\% lines full, \texttt{abd} excluded from \texttt{jacoco:check}), k6 load tests (3 runs, 50 VUs, 30\,s) with mean-of-means 23.01\,ms 95\%CI [-7.88; 53.91], max p95 173.13\,ms (threshold 200\,ms, 0\% errors) and OWASP ZAP baseline scan. Requirements follow ISO/IEC/IEEE 29148:2018 and all 15 INCOSE v4 characteristics were verified. SUS mean was 63.0 (SD 13.88, 95\%CI [53.07; 72.93], n=10). The system is versioned as v1.0.0 with Zenodo DOI.

\textbf{Keywords:} requirements engineering, SRS, Spring Boot, Angular, JWT, PostgreSQL, OWASP, ISO 29148, INCOSE, final course project.
